\section{Divisors and the lower-dimensional index input}
\label{sec:preliminaries}

We fix the divisor conventions needed for the exact pullback identities,
the rounded systems, and the index theorem on the reduced fibre.
All varieties are over $\C$.  Unless stated otherwise, they are integral
and normal.  In the global statements, all varieties are projective,
including the bases of the contractions.  The reduced fibres may be
reducible and nonnormal, and their normalizations may have several
connected components.  In the local curve arguments, the total spaces
are projective over the indicated curves.  After a curve is shrunk to
an open neighbourhood, neither it nor its total space need be projective
over $\C$.  A contraction is a surjective projective morphism $f$ with
$f_*\OO=\OO$.

Boundaries are effective rational divisors; crepant sub-boundaries may
have negative coefficients.  We use the usual discrepancy conventions
for lc, klt, and dlt pairs.  A reduced pair $(T,B_T)$ is slc if $T$ is
pure-dimensional, $S_2$, and nodal in codimension one, no component of
$B_T$ is contained in the conductor, $K_T+B_T$ is rational Cartier, and
the normalization with the conductor included in the boundary is lc.  In
particular, the conductor has coefficient one on each normalized branch.
The fibres to which we apply this definition will be proved to have the
required pure dimension and depth properties.

Canonical divisors are chosen using rational top forms, and we retain
those choices in exact divisor identities.  For rational divisors on a
normal variety, $D_1\sim_\Q D_2$ means that some positive integer multiple of
$D_1-D_2$ is the divisor of a rational function.  An equality of
rational divisors is coefficientwise for the chosen representatives.
In particular, $rD$ being an \emph{integral principal
divisor} means
\[
 rD=\Div(v)
\]
for an actual rational function $v$; it includes integrality of every
coefficient of $rD$.

For an integral Weil divisor $G$ on a normal variety $X$, the reflexive
divisorial sheaf is
\[
 \OO_X(G)(U)=\{v\in\C(X):\Div(v)|_U+G|_U\geq0\}\cup\{0\}.
\]
For a rational divisor $D$, its rounded system is the complete system of
$\OO_X(\lfloor D\rfloor)$.  A nonzero rational function $v$ is a section
exactly when
\[
 \Div(v)+D\geq0,
\]
because its valuations are integers.  This criterion does not require
$D$ or $\lfloor D\rfloor$ to be Cartier.  The field generated by a
nonempty system is the subfield of $\C(X)$ generated by all ratios of
nonzero sections.  It distinguishes the full embedded field from a
proper subfield of the same transcendence degree.

A subset of $\R$ satisfies the descending chain condition (DCC) if it
contains no infinite strictly decreasing sequence.  The ascending chain
condition (ACC) is defined by reversing the inequalities.  Both enter
the proof: threshold ACC gives a DCC set for discriminant coefficients,
and global ACC will reduce certain numerically trivial pairs to finite
coefficient sets.

The index input needed for the relative theorem is in fibre dimension at
most three.  We use the following established theorem: for every finite
set $I\subset[0,1]\cap\Q$, there is a positive integer $n(I)$ such that
every projective slc pair $(T,B_T)$ of dimension at most three with
coefficients in $I$ and $K_T+B_T\sim_\Q0$ satisfies
\begin{equation}\label{eq:lower-slc-index}
 n(I)(K_T+B_T)\sim 0.
\end{equation}
This includes normal pairs and the zero boundary
\cite[Corollary~1.6, arXiv version~1]{JiangLiu2021}.  We may enlarge
$n(I)$ to clear $I$.  The conclusion means that the log pluricanonical
sheaf $\OO_T(n(I)(K_T+B_T))$ is invertible and trivial on $T$ itself.  A
choice of generator is a global form whose restrictions to the
normalized components satisfy the conductor gluing.  This global
trivialization, rather than separate trivializations on the components,
is the input required by the residue argument.

We also need this conclusion when the coefficients initially lie in a
fixed rational DCC set $I\subset[0,1]$.  On each normalized irreducible
component, include the conductor with coefficient one.  The resulting
pair is projective lc with effective boundary coefficients in the DCC
set $I\cup\{1\}$, and its adjoint is numerically trivial.  Global ACC
\cite[Theorem~1.5]{HMX2014} places all its nonzero coefficients in a
fixed finite rational subset, depending only on $I$ and the dimension
bound.  The coefficients of $B_T$ therefore also lie in a fixed finite
set.  Applying~\eqref{eq:lower-slc-index} to the slc pair $T$ supplies one
uniform global trivialization, including the gluing.  The possibly larger
degree used later to restrict an ambient pluricanonical sheaf to $T$ may
depend on the ambient pair; it is kept separate from this uniform index.
